\subsection{Cache-Maintenance Blocks}
CPUID \texttt{CACHE\_TOPOLOGY} reports a cache-maintenance granule $G$. A cache-maintenance instruction computes and
translates its address-role EA without reading an operand value, then selects the physical byte interval
$[\lfloor A/G\rfloor G,\lfloor A/G\rfloor G+G)$ containing the translated (:term:base.terms.physical_address:) $A$. One instruction
operates on exactly that block. Because $G$ is at most 4096 bytes, the selected physical block does not cross a
16 KiB page boundary.

FLSHDCACHE, INVDCACHE, and WRBKDCACHE apply to every data or unified cache copy of the selected block in its
normal-memory coherence domain. FLSHDCACHE writes dirty bytes into normal-memory coherence and invalidates the
copies. INVDCACHE discards the copies without writeback. WRBKDCACHE writes dirty bytes into normal-memory coherence
and retains clean copies. Each operation completes the selected block before retirement.

INVICACHE invalidates the selected block in the executing (:term:base.terms.logical_processor:)\textquotesingle{}s instruction cache, decoded
instruction state, and instruction-prefetch state. SYNCCACHE first completes the data writeback for the block
throughout the coherence domain, then performs the same local instruction-side invalidation, and finally serializes
later instruction fetch so that it observes the synchronized bytes. To publish modified code to another logical
processor, software completes SYNCCACHE on the publisher, performs a release/acquire rendezvous, and executes
INVICACHE for the block on every receiving (:term:base.terms.logical_processor:) before branching to the modified bytes.

Address translation and permission checks precede a block\textquotesingle{}s cache effects. A fault leaves that instruction\textquotesingle{}s
selected block unchanged.
